\documentclass[10pt,a4paper]{amsart}
\usepackage[margin=19mm]{geometry}
\usepackage{amsmath,amssymb,mathtools}
\usepackage[scr=rsfs,cal=euler]{mathalfa}
\usepackage{xfrac}
\usepackage[unicode,hidelinks,bookmarks=false]{hyperref}
\newcommand{\ie}{i.e.\ }
\newcommand{\cf}{cf.\ }
\newcommand{\resp}{respectively\ }
\newcommand{\Subsection}{\subsection}
\providecommand{\nobreakdash}{\nobreak\hbox{-}}
\setlength{\emergencystretch}{2em}
\multlinegap0pt
\usepackage{amssymb}
\usepackage{float}
\usepackage{mathtools}
\usepackage{tikz}
\usepackage{xcolor}
\newtheorem{lem}{Lemma}
\newcommand{\NN}{\mathbb{N}}
\renewcommand{\d}{{\rm d}}
\title{Boundary dynamics, triple transitivity, and mixed identities in weakly hyperbolic groups}
\author{Ekaterina Rybak}
\date{Source: arXiv:2605.14159v3 (CC BY 4.0)}
\begin{document}
\maketitle

\noindent\emph{Source and licence.} Boundary dynamics, triple transitivity, and mixed identities in weakly hyperbolic groups. Version arXiv:2605.14159v3 (CC BY 4.0), distributed by arXiv under the Creative Commons Attribution 4.0 International licence, http://creativecommons.org/licenses/by/4.0/, as recorded in the arXiv licence metadata for this exact version and shown on the abstract page https://arxiv.org/abs/2605.14159v3. Full licence notice: \texttt{licenses/arxiv-2605.14159-CC-BY-4.0.txt}.\par\smallskip

\noindent\textit{Editorial scope.} Counted unit: the article's lem bounded (source0.tex lines 503--510). The article's complete original proof of it is delivered with it (source0.tex lines 512--560, one \texttt{proof} environment of 49 lines). No mathematics is written, repaired or re-derived: the statement and the proof are the article's own lines, in the article's own notation, and no retained line is substituted or re-flowed.  No further source-local context is needed: every cross-reference of every retained line resolves inside the two delivered spans.  Nothing was omitted from any retained span, so the package carries no omission record.  No retained line contains a citation, so the wrapper carries no bibliography and no bibliography entry.  No retained line contains a figure environment, an image-inclusion command or drawing code, so no image is delivered and the package declares no asset dependency.  Editorial formatting only, in addition: an AMS \texttt{amsart} preamble that restates the article's own declarations and macros used by the retained lines, namely the environments \texttt{lem} on separate counters, together with the standard packages those lines need.  The retained environments print in this document as Lemma 1 in order of appearance, whereas the article prints the same items with the numbering of its own full document.  The complete native source of the article as distributed in the arXiv e-print for this exact version is bundled unchanged as \texttt{sources/200560/source0.tex}.
\begin{lem}\label{gromov product is bounded}
Let $G$ be a group acting on $S$ by isometries. Let $(x_i)_{i\in \NN}$ be a sequence in $S$ converging to a point in $\partial S$. Fix a basepoint $x_0$ so that $\langle x_i \mid x_j \rangle_{x_0} \longrightarrow \infty$ as $i, j \longrightarrow \infty.$ and an arbitrary $g \in G$. Let $D=\d_S(x_0, gx_0)$. For a given $n \in \NN$ consider the geodesic quadrilateral $x_0, x_n, gx_n, gx_0$.
For $s \in [x_0, x_n]$, if 
\begin{equation}\label{Choice of s 2}
2D+4\delta < \d_S(x_0,s) <\langle x_n \mid  gx_n \rangle_{x_0} -2D -4\delta,
\end{equation}
then $\langle g^{-1}s \mid gs \rangle_s \leq 8\delta$.
\end{lem}

\begin{proof}
First, observe that for $s \in [x_0, x_n]$, if 
\begin{equation}\label{Choice of s}
D+2\delta < \d_S(x_0, s) < \langle x_n \mid gx_n \rangle_{x_0}- 2\delta,
\end{equation}
then $s$ is within $2\delta$ of some point $s' \in [gx_0,gx_n]$. Indeed, any geodesic quadrilateral in $S$ is $2\delta$-slim, so $s$ is within $2\delta$ of a point $s'$ on one of the sides $[x_0,gx_0]$, $[gx_0,gx_n]$, or $[gx_n,x_n]$.
If $s' \in [x_0, gx_0]$, then $$
\d_S(x_0,s) \leq \d_S(x_0,s')+\d_S(s',s) \leq D+2\delta,$$
which contradicts the first inequality in (\ref{Choice of s}).
If $s' \in [x_n, gx_n]$, then $$
\langle x_n \mid gx_n \rangle_{x_0} \leq \d_S(x_0, s') \leq \d_S(x_0,s)+2\delta,
$$ which contradicts the second inequality in (\ref{Choice of s}).

Therefore, if (\ref{Choice of s 2}) is satisfied, then there exists $s' \in [gx_0, gx_n]$ such that $\d_S(s, s') \leq 2\delta$. Consider the point $gs \in [gx_0, gx_n]$.
Note that, since $\d_S(x_0, gx_0) = D$ and $\d_S(s, s') \leq 2\delta$, by the triangle inequality, we obtain
$$
|\d_S(gx_0, s') - \d_S(x_0, s)| \leq D+2\delta.
$$
Thus, we have
$$
|\d_S(gx_0, s') - \d_S(gx_0,gs)| \leq D+2\delta.
$$
Since $gx_0, gs, s'$ are on the same geodesic $[gx_0, gx_n]$, we conclude that $L := \d_S(gs, s') \leq D+ 2\delta$. 
Consider the point $g^{-1}s' \in [x_0, x_n]$. We have $\d_S(s, g^{-1}s')=L$. Let $y$ be the other point on the geodesic $[x_0, x_n]$ at distance $L$ from $s$. Since $L\leq D+ 2\delta$, by (\ref{Choice of s 2}), we obtain that
$$
D+2\delta < \d_S(x_0, y) < \langle x_n \mid gx_n \rangle_{x_0}- 2\delta.
$$
Therefore, $y$ is within $2\delta$ of some point $y' \in [gx_0, gx_n]$.
Note that since $\d_S(s, s')\leq 2\delta$, by the triangle inequality
$|\d_S(y', s') - \d_S(y,s)| \leq 4\delta$.
Thus, we obtain
$$
L-4\delta \leq d_S(y', s') \leq L+4\delta.
$$
Since $\d_S(gs, s') =L$, and $y', gs, s'$ are all on the same geodesic, we get
$\d_S(y', gs) \leq 4\delta$.
We have $\d_S(g^{-1}s, g^{-1}s')=\d_S(s, s') \leq 2\delta$,
and $d_S(gs, y) \leq \d_S(gs, y')+\d_S(y, y') \leq 6\delta$.
So,
$$ 
\begin{aligned}
&\langle g^{-1}s, gs\rangle_s= \frac{1}{2}\left(\d_S(s, g^{-1}s)+\d_S(s, gs)-\d_S(g^{-1}s, gs)\right)\leq \\
%&\frac{1}{2}\left(\d_S(s, g^{-1}s')+\d_S(g^{-1}s',g^{-1}s)+ \d_S(s,y)+\d_S(y,gs)-\d_S(g^{-1}s', y)+\d_S(g^{-1}s, g^{-1}s')+\d_S(y,gs)\right) \\
&\leq \frac{1}{2}\left(\d_S(s, g^{-1}s')+2\delta
+ \d_S(s,y)+6\delta -\d_S(g^{-1}s', y)+8\delta)\right) \leq\\
&\leq \langle y \mid g^{-1}s'\rangle_s + 8\delta = 8\delta.
\end{aligned}
$$
\end{proof}

\end{document}
